\documentclass[11pt]{article}
\usepackage[T1]{fontenc}
\usepackage[margin=1in]{geometry}
\usepackage{amsmath,amssymb,amsthm,booktabs,array}
\usepackage{xurl}
\usepackage[hidelinks]{hyperref}
\newtheorem{theorem}{Theorem}
\newtheorem{proposition}[theorem]{Proposition}
\newtheorem{corollary}[theorem]{Corollary}
\newcommand{\F}{\mathbb F_3}
\title{Extending Finch's zero-run analysis:\\nonattainment and a corrected capture theorem}
\author{Haobo Ma\\\small ChronoAI Pte Ltd, Singapore
\and Wenlin Zhang\\\small National University of Singapore\\\small ChronoAI Pte Ltd, Singapore}
\date{November 19, 2026}
\begin{document}
\maketitle
\begin{abstract}
Extending Finch's analysis of the zero-run question in \emph{Three
Quadratic Recurrences Modulo 3}, arXiv:2610.34300v1, \S6, Remark~20,
p.~8, we answer universal attainment negatively: a closed ternary-state
orbit proves that the maximum at lag $61$ is ten gaps rather than the
upper bound fifteen. The ten-gap starts are exactly $2780+8946m$
for $m\ge0$, and the least preperiod is $739$; Harder's earlier period
value $8946$ is recovered. We also correct Finch's capture-offset
Conjecture~23 (\S7, p.~9): whenever the eventual period is four,
the offset is six for every lag $k\ge3$, and nine at $k=2$.
\end{abstract}

\section{The question and the sharp counterexample}
Fix $k\ge2$. The positive-index sequence considered here is uniquely
specified in $\F$ by
\begin{equation}\label{eq:recurrence}
 x_1=\cdots=x_{k-1}=0,\qquad x_k=1,\qquad
 x_n=x_{n-1}+x_{n-k}^{\,2}\quad(n\ge k+1).
\end{equation}
Represent its values by $0,1,2$. Following Finch's Definition~15
\cite[\S6, pp.~6--7]{finch2026}, a run with $g\ge1$ gaps starts at
$r\ge k+1$ and satisfies
\begin{equation}\label{eq:run}
 x_{r+i}=0\quad\Longleftrightarrow\quad i\equiv0\pmod4
 \qquad(0\le i\le4g).
\end{equation}
It has $g+1$ zeros and need not be a maximal run.

Finch's Theorem~16 bounds $g$ by $\lfloor(k-1)/4\rfloor$ for
$k\equiv0,1\pmod4$. Remark~20 asks about attaining that bound
for every such parameter. Our result concerns this universal reading
of the question, not the validity of the upper bound. Earlier period
computations by Harder, reproduced in Finch's note
\cite[\S0.2, p.~2]{finch2002}, include $N_A(61)=8946$.
Period length alone does not decide whether a specified zero pattern occurs;
the complete orbit and the pattern test below do.

\begin{theorem}\label{thm:main}
For the sequence \eqref{eq:recurrence} with $k=61$, the maximum number
of gaps in a run \eqref{eq:run} is exactly ten. In particular, the
bound $\lfloor(61-1)/4\rfloor=15$ is not attained. The least eventual
period is $8946$, and the least $\mu\ge0$ such that
$x_{n+8946}=x_n$ for all $n>\mu$ is $739$.
The ten-gap runs start exactly at $2780+8946m$, for integers $m\ge0$.
\end{theorem}

\section{A finite certificate for the entire sequence}
For $j\ge0$ encode the block ending at $j+61$ as
\[
 S_j=\sum_{i=0}^{60}x_{j+61-i}3^i.
\]
Thus $S_0=1$, and the newest term is the least significant digit.
Writing $M=3^{60}$, define, for $0\le s<3^{61}$,
\begin{equation}\label{eq:transition}
 T(s)=3(s\bmod M)+\left(\bigl((s\bmod3)+\lfloor s/M\rfloor^2\bigr)\bmod3\right).
\end{equation}
In the parentheses the whole sum is reduced modulo $3$.
Indeed $s\bmod3=x_{j+61}$ and $\lfloor s/M\rfloor=x_{j+1}$ when
$s=S_j$. Multiplication by three shifts the retained sixty digits
one place; the new low digit is $x_{j+62}$. Hence
$S_{j+1}=T(S_j)$ for every $j\ge0$.

The exact-integer calculation in \texttt{verify.py} gives
\begin{align}
 &S_0,\ldots,S_{9684}\text{ are pairwise distinct},\label{eq:distinct}\\
 &T(S_{9684})=S_{739}=92196811829875935758365998383.\label{eq:closing}
\end{align}
Every successor in the list is checked as well. Thus these $9685$
states form a forward-invariant set containing $S_0$. Induction
puts every future state in the set; a finite-prefix observation
has therefore become a statement about the entire infinite sequence.

For $1\le h\le15$, let $P_h(s)$ mean
\[
 \left(\left\lfloor s/3^i\right\rfloor\bmod3=0\right)
 \Longleftrightarrow i\equiv0\pmod4
 \qquad(0\le i\le4h).
\]
The same certificate gives
\begin{center}
\begin{tabular}{@{}lrrr@{}}
\toprule
Mask & $P_{10}$ & $P_{11}$ & $P_{15}$\\
\midrule
Number of states among $S_0,\ldots,S_{9684}$ & 1 & 0 & 0\\
\bottomrule
\end{tabular}
\end{center}
These counts exhaust the finite invariant set, rather than a
chosen number of scalar terms. The sole state satisfying $P_{10}$
is $S_{2759}$. The transition, distinctness, and mask equalities above
are the finite premises of the proof; the complete calculation and
its integer-valued output are in \texttt{verify.py} and
\texttt{certificate.json}.

\begin{proof}[Proof of Theorem~\ref{thm:main}]
Suppose a run with at least eleven gaps starts at $r\ge62$.
Its first eleven gaps cover $r,\ldots,r+44$.
The lowest $45$ digits of $S_{r-17}$ are, in reverse order,
$x_r,\ldots,x_{r+44}$. Since $44\equiv0\pmod4$, reversal preserves
the prescribed zero positions, so $P_{11}(S_{r-17})$ holds.
Here $r-17\ge45$, so all these digits correspond to positive indices.
The invariant-set calculation excludes that mask for every reachable
state. This proves $g\le10$.

For equality, direct substitution in \eqref{eq:recurrence} gives
\[
 (x_{2780},\ldots,x_{2820})=(0,1,2,2)^{10}0,
\]
where the exponent denotes concatenation. All forty-one terms are
checked, and $2780\ge62$, so this is an admissible ten-gap run.

Finally, a deterministic orbit whose first repetition is
\eqref{eq:closing}, with \eqref{eq:distinct}, has least preperiod
$739$ and least cycle length $9685-739=8946$.
For completeness, eventual scalar periodicity with period $p$ after
$\mu$ implies $S_{j+p}=S_j$ for $j\ge\mu$, by the block encoding.
Conversely, periodic states for $j\ge\mu$ imply
$x_{n+p}=x_n$ for $n>\mu$, by their most significant digits.
Thus the least eventual periods agree, as do their least preperiods.

A ten-gap run at $r\ge62$ is equivalent to $P_{10}(S_{r-21})$:
the same reversal argument uses forty-one digits. The sole such
state occurs at $j=2759$ in the cycle starting at $739$, and hence
exactly at $j=2759+8946m$ for $m\ge0$. This gives the asserted starts.
\end{proof}

The scalar recurrence provides a second implementation of the same
finite premises. Its comparison ranges are recorded in
\texttt{certificate.json}; the invariant-set argument supplies the
infinite scope.

\section{The corrected capture theorem}
For the same recurrence, Finch's Definitions~21--22 and Conjecture~23
\cite[\S7, p.~9]{finch2026} use $G=\lceil(k-1)/4\rceil$ and the
least endpoint $z$ of an admissible $G$-gap run. When the eventual
period is four, $\mu_A(k)$ is the least nonnegative integer with
$x_{n+4}=x_n$ for every $n>\mu_A(k)$. The proposed identity is
$z-\mu_A(k)=6$ for $k\equiv2,3\pmod4$.
The four-cycle patterns go back to the discussion recorded by Miller
\cite{finchMO}; Finch's Propositions~12--13 \cite[\S4, p.~5]{finch2026}
classify the cycles and give the capture criterion. The correction
below uses those results; it does not determine which lags enter a
four-cycle.

\begin{theorem}\label{thm:capture}
Suppose $k\ge2$ and the least eventual period of
\eqref{eq:recurrence} is four. Then
\[
z(k)-\mu_A(k)=
\begin{cases}9,&k=2,\\6,&k\ge3.\end{cases}
\]
In particular, the restriction of Conjecture~23 to $k\ge3$ is true.
For $k\ge3$, necessarily $k\equiv2,3\pmod4$, and the first
capturing run starts at
\[
s=\begin{cases}
\mu_A(k)-k+4,&k\equiv2\pmod4,\\
\mu_A(k)-k+5,&k\equiv3\pmod4.
\end{cases}
\]
The triple $(x_{\mu_A(k)},x_{\mu_A(k)+1},x_{\mu_A(k)+2})$
is respectively $(2,2,0)$ or $(1,2,0)$.
\end{theorem}
\begin{proof}
First take $k=2$.
The repeating block $(0,1,1,2)$ gives the initial values and satisfies
the recurrence: at residues $n\equiv0,1,2,3\pmod4$, the respective
identities in $\F$ are
$2=1+1^2$, $0=2+1^2$, $1=0+2^2$, and $1=1+0^2$.
Uniqueness of the recurrence proves the description for every $n\ge1$.
The zeros are exactly $n\equiv1\pmod4$, so the least period is four
and the least preperiod is zero. Here $G=1$ and an admissible run
must start at $r\ge3$. The first possible zero is $r=5$, followed
by the next at $9$, with nonzero terms between. Hence $z=9$.

Now let $k\ge3$ and write $\mu=\mu_A(k)$ and
$b_n=[x_n\ne0]$. Let $(c_n)_{n\in\mathbb Z}$ be the
four-periodic sequence agreeing with the tail of $x$, and put
$\beta_n=[c_n\ne0]$. Shifting an index by a sufficiently large
multiple of four shows that $c$ satisfies the recurrence at every
integer index.

For completeness, Finch's Proposition~12 gives the following
classification by an elementary increment count. The increments
$c_n-c_{n-1}=\beta_{n-k}$ around a period sum to zero in $\F$,
so the number of nonzero entries is divisible by three. It is not
zero, since the least period is four; hence there is exactly one
zero. Starting there, the block is $(0,1,a,2)$ with $a\in\{1,2\}$:
the first and last steps force the entries one and two. The sole
zero increment is at offset two if $a=1$, and at offset three if
$a=2$. Its delayed argument is the sole zero, so necessarily
$k\equiv2\pmod4$ when $a=1$, and $k\equiv3\pmod4$ when $a=2$.
This recovers the cycle patterns credited to Pinch and Finch above.

The initial conditions give $x_k=\cdots=x_{2k-1}=1$. The longest
string of ones in $c$ has length two when $a=1$ and one when $a=2$.
Consequently
\[
\mu\ge\begin{cases}2k-3,&a=1,\\2k-2,&a=2.\end{cases}
\]
In the first case $k\ge6$, and in the second $k\ge3$; thus
$\mu\ge k+1$. Minimality of $\mu$ gives $x_\mu\ne c_\mu$,
and $x_n=c_n$ for all $n>\mu$.

Set $L=\mu+2-k$, so $L\ge1$. Comparing the recurrences at $n+k$
for $n\ge L$ gives
\begin{equation}\label{eq:indicator}
b_n=\beta_n\quad(n\ge L),\qquad b_{L-1}\ne\beta_{L-1}.
\end{equation}
The second assertion follows by comparing the recurrences at
$\mu+1$: equality of these indicators would force
$x_\mu=c_\mu$. All indices here are positive, since the displayed
lower bounds on $\mu$ give $L-1\ge1$.
More explicitly, for $n\ge L$,
\[
b_n=x_{n+k}-x_{n+k-1}=c_{n+k}-c_{n+k-1}=\beta_n.
\]

In particular, $b_\mu=\beta_\mu$; hence $x_\mu,c_\mu$ are distinct
nonzero residues. The discrepancy cannot be at the first nonzero
position after a zero of $c$: then $c_\mu=1$, $x_\mu=2$, and
$x_{\mu-1}=0$ by \eqref{eq:indicator}, giving the impossible step
$2-0=2$. It cannot be at the last nonzero position either: then
$x_\mu=1$, $x_{\mu+1}=0$, again giving a step of two.
Every step used has index at least $k+1$, so its value must be zero
or one. Thus $\mu$ is exactly the middle position of the block,
two places after a zero of $c$.
Consequently the last-disagreement triples are $(2,2,0)$ when
$a=1$ and $(1,2,0)$ when $a=2$. The recurrence at $\mu+1$ gives
\[
b_{L-1}=c_{\mu+1}-x_\mu=
\begin{cases}0,&k\equiv2\pmod4,\\1,&k\equiv3\pmod4.\end{cases}
\]

Put $s=\mu+6-4G$. For $k=4q+2$ we have $4G=k+2$,
$s=L+2$, and $\mu\ge2k-3$ gives $s\ge k+1$.
For $k=4q+3$ we have $4G=k+1$, $s=L+3$, and
$\mu\ge2k-2$ gives $s\ge k+3$. These are exactly the two
formulas for $s$ in the theorem. In each case $s$ is
a zero position of $c$. Equation~\eqref{eq:indicator} therefore
gives a $G$-gap run from $s$ to $\mu+6$, so $z\le\mu+6$.

For minimality, recall the forward implication of the capture
criterion: a $G$-gap run ending at $e$ determines $x_e=0$ and
$k$ consecutive indicators ending at $e$, since $4G+1\ge k$.
The data $(x_e;b_{e-k+1},\ldots,b_e)$ coincide with the
four-cycle having a zero at $e$; they determine $x_{e+1}$,
then its indicator, and inductively every later term. Thus $x$ agrees
with $c$ from $e$ onward. No such endpoint can be at or before
$\mu$. After $\mu$ the first two zeros are $\mu+2$ and $\mu+6$.
A run ending at $\mu+2$ would start at $s-4$. If $k\equiv2$,
this start is $L-2$, but the interior position $L-1$ is an extra
zero by \eqref{eq:indicator}. If $k\equiv3$, the proposed start is
$L-1$, which is nonzero by the same equation. Both are impossible.
Hence $z=\mu+6$.
\end{proof}

\begin{corollary}\label{cor:neighborhood}
Under the hypotheses of Theorem~\ref{thm:capture}, put $s=z-4G$.
For $k\ge3$ the neighborhood conditions in Finch's Remark~24 hold:
$b_{s-4}=1$ when $k\equiv2\pmod4$, and
$b_{s-5}=0$, $b_{s-1}=1$ when $k\equiv3\pmod4$.
At $k=2$ the proposed condition fails: $s=5$ and $b_{s-4}=b_1=0$.
\end{corollary}
\begin{proof}
Use $L,\mu,a$ from the proof. If $a=1$, then $k\ge6$ and
$\mu-2\ge L$, so $x_{\mu-2}=0$, whereas $x_{\mu-1}\ne0$.
The step condition gives $x_{\mu-1}=1$. If $a=2$, then
$x_\mu=1$ and $x_{\mu-1}\ne0$; the step into $x_\mu$ also
forces $x_{\mu-1}=1$. The recurrence at $\mu$ now gives
$b_{\mu-k}=1$ for $a=1$ and $b_{\mu-k}=0$ for $a=2$.
These are $b_{s-4}$ and $b_{s-5}$ respectively.
In the latter case $s-1=L+2$ is nonzero by
\eqref{eq:indicator}. The $k=2$ value follows from its explicit cycle.
\end{proof}

\section{Coverage and formal statements}
\begin{center}\small
\begin{tabular}{@{}>{\raggedright\arraybackslash}p{0.28\linewidth}>{\raggedright\arraybackslash}p{0.65\linewidth}@{}}
\toprule
Item in Finch \cite{finch2026} & Coverage\\
\midrule
Definition~15; Remark~20 & Answered: no, for universal attainment.
The exact maximum at $k=61$ is ten (Theorem~\ref{thm:main}).\\
Theorem~16; Lemma~17 & Prior proof cited: Finch (2026).
The bound and local block description remain valid.\\
Propositions~12--13; Remark~19 & Classification reproved; prior proof:
Finch (2026). Capture criterion cited; both are used above.\\
Example~18 & Prior computation cited: Finch (2026); the sixty
terms agree with the recurrence.\\
Remark~20, six listed lags & Not covered: their individual attainment
claims are distinct from universal attainment.\\
Conjecture~23 & Refuted-and-corrected: offset nine at $k=2$;
offset six at every $k\ge3$ with period four
(Theorem~\ref{thm:capture}).\\
Remark~24 & Neighborhood conditions proved for $k\ge3$;
the boundary condition fails at $k=2$
(Corollary~\ref{cor:neighborhood}).\\
Conjectures~5, 10, 14 & Not covered. The capture theorem is
conditional on period four, not a proof of either infinite family
in Conjecture~14.\\
\bottomrule
\end{tabular}
\end{center}

\subsection*{Supporting finite checks}
The program \texttt{verify.py} closes all canonical state orbits for
$2\le k\le41$ and checks the capture and neighborhood values in
the seventeen period-four cases. The separate program
\texttt{verify\_capture.py} enumerates all $81$ four-letter words
for each lag residue and searches the forty parameters
$2\le k\le80$ with $k\equiv2,3\pmod4$, up to index $10^6$
per parameter. It finds thirty capturing runs and verifies their
cycle states, start indices, last-disagreement triples, and local
indicators. For $k=18,23,35,42,50,51,55,63,74,78$ it finds no
capture within that bound; this gives no classification of their
eventual periods. The exact outputs are \texttt{certificate.json}
and \texttt{capture-checks.json}. These examples supplement the
written all-parameter proofs.

\newpage
\subsection*{Formal coverage}
The declaration \texttt{result} in the pinned
\href{https://github.com/the-omega-institute/trureturing/blob/0d064f0d28d301cd8a1d1b09173035a9f2fd6c02/D5/S1/Recurrence/FinchZeroRunsModThree.lean}{\texttt{FinchZeroRunsModThree.lean}}
proves the negation of universal attainment, using the $k=61$
orbit and the excluded eleven-gap mask. The explicit ten-gap block
and the least-period and preperiod conclusions above have the stated
finite-check and written-proof coverage.
The declaration \texttt{result} in the pinned
\href{https://github.com/the-omega-institute/trureturing/blob/0d064f0d28d301cd8a1d1b09173035a9f2fd6c02/D5/S0/Certificates/FinchConjecture23CaptureOffsetRefutation.lean}{capture-offset Lean source}
proves the negation of the capture-offset conjecture using the $k=2$
cycle, least capture endpoint, and least preperiod in
Theorem~\ref{thm:capture}. The $k\ge3$ capture correction and
Corollary~\ref{cor:neighborhood} have the written proofs above,
additional to the conclusions of these two Lean declarations.

\begin{thebibliography}{9}
\bibitem{finch2026} S. Finch,
\emph{Three Quadratic Recurrences Modulo 3: What We Can Prove and What We Can Guess},
arXiv:2610.34300v1 (2026), \S\S6--7, pp.~6--10.
\url{https://arxiv.org/abs/2610.34300v1}.
\bibitem{finch2002} S. R. Finch,
\emph{Periodicity in Sequences Mod 3}, September 5, 2002;
\S0.2 records B. Harder's computations from 2000.
\url{https://www.mit.edu/~sfinch/seqmod3.pdf}.
\bibitem{finchMO} V. S. Miller,
\emph{Finch's sequence over $\mathbb F_3$},
MathOverflow question 50678, December 29, 2010; the question records
R. Pinch's explanation of the four-periodic patterns.
\url{https://mathoverflow.net/questions/50678/finchs-sequence-over-mathbbf-3}.
\end{thebibliography}
\bigskip
\noindent\small\copyright\ 2026 Haobo Ma and Wenlin Zhang.
Licensed under Creative Commons Attribution 4.0
International (CC BY 4.0).
\end{document}
